\section{Travaux connexes}
\label{sec:related}

\paragraph{Apprentissage continu embarqué.} Les méthodes de CL se répartissent classiquement en approches
par régularisation, par rejeu et par architecture~\cite{DeLange2021Survey}. EWC~\cite{Kirkpatrick2017EWC}
appartient à la première famille : une pénalité quadratique pondérée par l'information de Fisher protège
les paramètres importants pour les tâches passées, sans stocker de données --- une propriété précieuse
sous contrainte mémoire. Sur MCU, \gls{TinyOL}~\cite{Ren2021TinyOL} a montré la faisabilité de l'apprentissage
en ligne d'une couche additionnelle, et \gls{HDC}~\cite{Benatti2019HDC} celle d'un apprentissage non neuronal à
très basse consommation. LifeLearner~\cite{Kwon2023LifeLearner} et les plateformes à rejeu latent
quantifié~\cite{Ravaglia2021QLRCL} adressent le compromis mémoire/précision sur des cibles embarquées. Les
panoramas \gls{TinyML} récents~\cite{Capogrosso2023TinyML} soulignent que l'entraînement \emph{sur l'appareil}
reste sous-exploré comparé à l'inférence seule, et les travaux appliqués à la maintenance
prédictive~\cite{Hurtado2023CLPdM} restent majoritairement évalués hors cible.

\paragraph{Quantifier : trois degrés de liberté, rarement croisés.} On résume souvent la quantification à
un choix de \emph{format} --- FP32 vers INT8, ou virgule fixe Q15. Deux autres degrés la commandent
pourtant tout autant. Le \emph{moment} d'abord~\cite{Jacob2018,Krishnamoorthi2018} : quantifier après l'entraînement, en PTQ, est immédiat mais
sensible à la calibration de l'échelle, tandis que quantifier pendant, en QAT, coûte un réentraînement et
suppose un accès au pipeline d'apprentissage, ce dont on ne dispose pas toujours en contexte industriel.
La \emph{profondeur} ensuite, avec la granularité qui l'accompagne : descendre sous 8 bits ne vaut que si
l'échelle est estimée assez finement, par canal plutôt que par tenseur, et que si le stockage suit
réellement le format annoncé. Ces trois axes sont le plus souvent étudiés séparément ; nous les faisons
varier sur \emph{une même tête, un même jeu et une même carte}, ce qui permet de les comparer plutôt que
de les juxtaposer.

\paragraph{Ce qui est rarement vérifié : les trois promesses, séparément.} La quantification est
habituellement justifiée par un argument tripartite --- moins de mémoire, moins de temps, moins d'énergie
--- hérité de cibles disposant d'unités entières dédiées. Sur un cœur généraliste doté d'une \gls{FPU}, rien ne
garantit que les trois termes se comportent de la même façon, et les rapporter conjointement suppose une
instrumentation que peu de travaux mettent en place : mesurer la mémoire \emph{totale} du système et non
la seule empreinte des poids, chronométrer le noyau poste par poste, et intégrer un courant réel à la
sonde. Notre contribution empirique est de traiter ces trois promesses comme trois questions distinctes,
mesurées sur la même carte : nous montrons que la première est tenue mais porte sur un poste minoritaire,
et que les deux autres ne le sont pas. Nous établissons au passage l'effondrement de la PTQ naïve d'une
tête de décision, en isolons la cause plutôt que de la masquer par un QAT en amont, et en démontrons la
récupération par un noyau calibré.
